\section*{Capítulo \ref{ch:g5:natural-hazards} --- As ameaças naturais}

\begin{solution}{exo:g5:natural-hazards:1}
(a) Só uma \omterm{def:g5:natural-hazards:hazard}{ameaça natural}: ninguém sai ferido. (b) Um \omterm{def:g5:natural-hazards:disaster}{desastre natural}:
casas são danificadas. (c) Só uma \omterm{def:g5:natural-hazards:hazard}{ameaça natural}. (d) Um \omterm{def:g5:natural-hazards:disaster}{desastre natural}.
\end{solution}

\begin{solution}{exo:g5:natural-hazards:2}
Do ar: uma onda de calor, uma \omterm{def:g5:natural-hazards:storm}{tempestade}, uma \omterm{def:g5:natural-hazards:drought}{seca}. Da água: uma
\omterm{def:g5:water-shapes-the-land:flood}{inundação}. Do chão: um \omterm{def:g5:natural-hazards:landslide}{deslizamento de terra}, uma \omterm{def:g4:volcanoes-earthquakes-first-look:eruption}{erupção} vulcânica, um
\omterm{def:g4:volcanoes-earthquakes-first-look:earthquake}{terremoto}.
\end{solution}

\begin{solution}{exo:g5:natural-hazards:3}
Uma \omterm{def:g5:natural-hazards:drought}{seca} é um período longo, de meses ou de anos, com muito menos \omterm{def:g4:clouds-rain-snow:precipitation}{chuva} que
o normal. Ela pode deixar os agricultores sem colheita, porque as
plantações murcham e morrem, e pode deixar as cidades com falta de água
potável, porque os \omterm{def:g1:landscapes-around-us:water}{rios} e os \omterm{def:g1:landscapes-around-us:water}{lagos} encolhem.
\end{solution}

\begin{solution}{exo:g5:natural-hazards:4}
Por exemplo: guardar ou amarrar tudo o que o vento possa levar (cadeiras
de jardim, brinquedos, latas de lixo); conferir as lanternas e deixar
pilhas de reserva e água à mão; ficar dentro de casa, longe das janelas,
durante a \omterm{def:g5:natural-hazards:storm}{tempestade}; continuar escutando as notícias no rádio.
\end{solution}

\begin{solution}{exo:g5:natural-hazards:5}
(a) Um furacão. (b) Um tufão. (c) Um ciclone tropical.
\end{solution}

\begin{solution}{exo:g5:natural-hazards:6}
Verde, amarelo, laranja, vermelho. No nível vermelho, o perigo é grande:
aja na mesma hora e siga as instruções dos responsáveis.
\end{solution}

\begin{solution}{exo:g5:natural-hazards:7}
Uma parte da encosta da direita se soltou e escorregou, deixando uma
cicatriz nua. A sua \omterm{def:g3:stones-pebbles-sand:sand}{lama} se espalhou por todo o fundo do \omterm{def:g1:landscapes-around-us:landform}{vale} e cobriu a
estrada. Ela também bloqueou o \omterm{def:g1:landscapes-around-us:water}{rio} como uma barragem: a água do \omterm{def:g1:landscapes-around-us:water}{rio}, que
não consegue mais escoar, se acumulou atrás dela e inundou a mata \omterm{def:g1:landscapes-around-us:water}{rio}
acima.
\end{solution}

\begin{solution}{exo:g5:natural-hazards:8}
Lena tem razão. Não podemos impedir a \omterm{def:g4:clouds-rain-snow:precipitation}{chuva} nem o \omterm{def:g1:landscapes-around-us:water}{rio}, mas podemos manter
as pessoas longe do perigo. Por exemplo: construir as casas num terreno
mais alto, e não na \omterm{def:g5:water-shapes-the-land:flood}{planície de inundação}; avisar as pessoas a tempo com
alertas de \omterm{def:g5:water-shapes-the-land:flood}{inundação}, para que elas subam para um terreno mais alto;
ensinar a todos que nunca se deve atravessar a água de uma enchente a pé
ou de carro.
\end{solution}

\begin{solution}{exo:g5:natural-hazards:9}
(a) $4 \times 60 = 240$ milímetros. (b) Um quarto de 60 é 15; em quatro
meses, $4 \times 15 = 60$ milímetros. (c) Faltaram $240 - 60 = 180$
milímetros: três quartos da \omterm{def:g4:clouds-rain-snow:precipitation}{chuva} normal.
\end{solution}

\begin{solution}{exo:g5:natural-hazards:10}
(a) $1\,200 \div 4 = 300$ casas. (b) Um terço de 300 é 100; então
$300 - 100 = 200$ casas continuam na \omterm{def:g5:water-shapes-the-land:flood}{planície de inundação}.
\end{solution}

\begin{solution}{exo:g5:natural-hazards:11}
Motivos possíveis: a primeira vila foi avisada a tempo e as pessoas se
prepararam (janelas fechadas, objetos guardados), e a segunda não; as
casas da primeira são construídas de forma sólida, e as da segunda são
frágeis; a segunda vila fica num \omterm{def:g1:landscapes-around-us:water}{litoral} baixo ou numa \omterm{def:g5:water-shapes-the-land:flood}{planície de
inundação}, onde a \omterm{def:g5:natural-hazards:storm}{tempestade} também trouxe uma \omterm{def:g5:water-shapes-the-land:flood}{inundação}, enquanto a
primeira fica num terreno mais alto; na segunda vila, árvores perto das
casas caíram em cima delas. A mesma ameaça não causa os mesmos danos em
todo lugar: isso depende de onde e de como as pessoas vivem, e de quão bem
elas estão preparadas.
\end{solution}
